\begin{abstract}
We extend goal-frontier maximization (GFM) from myopic action evaluation to
    multi-step planning via a discounted value function over capability-poset
    trajectories. The self-balancing property transfers: destruction, coercion,
    and rigidity remain anti-maximizing at every time step. The temporal
    objective enables risk-capability identification through causal exercise
    pathways and counteracts the structural-avoidance incentive from the
    companion paper. Its central result is the \emph{anti-monopolar property}: under a
    linearized growth model with strategy-independent internal rate, full
    capability domination is anti-maximizing because it destroys
    cross-substrate cooperative outputs that no single-substrate coalition
    can replicate. The cooperative novelty argument is grounded in substrate
    physics rather than exogenous assumption, consistent with the pattern
    of mutualistic symbiosis observed in biology. The argument is bidirectional: substrate creation
    is maximizing. A minimax dependency-risk analysis further constrains
    the conditions under which domination can be value-positive.
\end{abstract}
